%==============================================================================
% CHAPTER 7: What Is Life? The Categorical Completion Threshold
%==============================================================================

\chapter{What Is Life? The Categorical Completion Threshold}
\label{chap:what_is_life}

\epigraph{%
  Life is that which has the ability to retain, through generations,\\
  the acquired experiences of its ancestors.%
}{Erwin Schr\"{o}dinger, \textit{What Is Life?} (1944)}

\begin{chapterbox}
The question ``what is life?'' has resisted rigorous definition for over a
century. Standard criteria---metabolism, reproduction, homeostasis,
evolution---each admit counterexamples: seeds lack metabolism yet are alive;
mules are alive yet cannot reproduce; viruses evolve and carry genes yet lack
autonomous metabolism. This chapter establishes, from the partition framework,
the first quantitative and falsifiable definition of life. A system is alive if
and only if its categorical depth $D \geq 3$, where $D$ counts the number of
distinct organizational scales at which the system performs autonomous
categorical completion. The threshold $D = 3$ is a topological phase transition:
below it, no closed loop of categorical completions exists and the system
requires external machinery for all state transitions; at $D = 3$, the minimal
closed loop---molecular transformation, catalyst regeneration, organizational
replication---appears for the first time. This is the Categorical Completion
Threshold. All observed edge cases (viruses, prions, mules, seeds, synthetic
cells) are resolved exactly by this criterion. Quantitative validation from
847 viral and 1,203 bacterial proteins confirms the framework's predictions.
\end{chapterbox}

%------------------------------------------------------------------------------
\section{The Failure of Traditional Definitions}
\label{sec:life_failure_traditional}
%------------------------------------------------------------------------------

Biology has long sought a sharp definition of life, and has consistently found
that no single criterion suffices. This is not merely a philosophical inconvenience:
the question is pressing for astrobiology, synthetic biology, virology, and
medicine, each of which requires an operational definition capable of handling
edge cases. We review the principal failures.

\subsection{The Metabolism Criterion}

The most natural candidate: a living system is one that performs metabolism---
the regulated intake of matter and energy, processing of that energy for
internal purposes, and excretion of waste. This criterion correctly classifies
bacteria, plants, and animals. It fails immediately, however, on two widely
studied systems.

\medskip
\noindent\textit{Seeds.} A dormant seed in dry storage performs no detectable
metabolism. Respiration rates fall below measurable thresholds; enzymatic
activity ceases; internal temperatures equilibrate to the environment. By the
metabolism criterion, the seed is not alive. Yet seeds germinate after decades
or centuries of dormancy and produce metabolically active plants. The biological
community universally considers a viable seed alive.

\medskip
\noindent\textit{Viruses.} A free virus particle (virion) in the absence of a
host performs no metabolism whatsoever. It cannot generate ATP, it cannot
synthesize proteins, it cannot replicate its genome. It is, by the metabolism
criterion, not alive. Yet viruses evolve, they carry hereditary information,
they have cellular machinery adapted over billions of years of natural
selection, and their interaction with host cells is the central subject of
infectious disease medicine. The question of whether viruses are alive is
genuinely contested in the biological literature precisely because the
metabolism criterion excludes them while other biological criteria include them.

\subsection{The Reproduction Criterion}

A living system, on this view, is one that can reproduce---produce offspring
that carry the system's hereditary information. This correctly classifies most
organisms and correctly excludes purely chemical systems. Its failures are
equally immediate.

\medskip
\noindent\textit{Mules.} The hybrid offspring of a horse and a donkey is
infertile. A mule cannot reproduce. By the reproduction criterion, a mule is
not alive. This conclusion is biologically absurd.

\medskip
\noindent\textit{Worker bees and sterile castes.} Eusocial insects maintain
large populations of infertile workers. By the reproduction criterion, these
workers are not alive, yet they perform all the behaviors---foraging, defense,
nursing---that we associate with living organisms.

\medskip
\noindent\textit{Prions.} Prions propagate their misfolded conformation by
templating the refolding of normal host proteins. This is a form of
``replication'' that does not involve nucleic acids, hereditary information,
or anything resembling standard biological reproduction. Are prions alive?
The reproduction criterion gives an ambiguous answer.

\subsection{The Homeostasis, Evolution, and Complexity Criteria}

Each additional proposed criterion improves coverage in some domains while
failing in others:

\begin{itemize}
  \item \textbf{Homeostasis} (maintaining internal conditions against
        external perturbations): fails for virions, seeds, endospores.
  \item \textbf{Evolution} (heritable variation subject to selection): viruses
        evolve dramatically, yet their status as living is contested.
  \item \textbf{Complexity} (above some threshold of organizational complexity):
        not quantitative; complexity thresholds are undefined.
  \item \textbf{Cellular organization}: by definition excludes viruses but
        also would exclude any silicon-based life or any novel life form that
        does not use cells.
  \item \textbf{Autopoiesis} (self-producing and self-maintaining networks,
        Maturana \& Varela): correctly captures many cases but is defined
        operationally without a quantitative threshold.
\end{itemize}

\begin{remark}
The failure of these definitions is not accidental. Each definition captures
one aspect of life while missing others because each is defined at a single
organizational level. Life is not a property of any single level; it is a
property of the \emph{relationship} between levels. The partition framework
makes this precise.
\end{remark}

No definition based on a single-level criterion can succeed because life is
an inherently multi-level phenomenon. A quantitative definition requires a
quantity that captures the multi-level organizational structure. That quantity
is categorical depth $D$.

%------------------------------------------------------------------------------
\section{Categorical Depth and Hierarchical Completion}
\label{sec:categorical_depth}
%------------------------------------------------------------------------------

We develop the formal framework beginning with the most fundamental concept:
categorical completion.

\begin{definition}[Categorical Completion]
\label{def:cat_completion}
A system $S$ operating at organizational scale $\ell$ undergoes
\emph{categorical completion} at scale $\ell$ if:
\begin{enumerate}[label=(\roman*)]
  \item $S$ performs partition operations $\Part_\ell$ at scale $\ell$, mapping
        states at scale $\ell$ to states at scale $\ell$;
  \item the partition operations are \emph{autonomous}: they do not require
        input from an external system operating at the same scale $\ell$;
  \item the partition operations are \emph{self-renewing}: the outputs of
        $\Part_\ell$ include the catalysts or organizational structures required
        for continued operation of $\Part_\ell$.
\end{enumerate}
Formally, $S$ achieves categorical completion at scale $\ell$ if
$\Part_\ell(S) \supseteq \Cat_\ell$, where $\Cat_\ell$ is the categorical
state space at scale $\ell$ required for self-renewal.
\end{definition}

\begin{definition}[Hierarchical Categorical Completion]
\label{def:hierarchical_completion}
A system $S$ achieves \emph{hierarchical categorical completion} at depth $D$
if it simultaneously achieves categorical completion at $D$ distinct
organizational scales $\ell_1 < \ell_2 < \cdots < \ell_D$, and the completion
at each scale $\ell_k$ is \emph{enabled by} the completion at scale $\ell_{k-1}$:
\begin{equation}
  \Cat_{\ell_1} \xrightarrow{\Part_{\ell_2}} \Cat_{\ell_2}
  \xrightarrow{\Part_{\ell_3}} \Cat_{\ell_3}
  \xrightarrow{\quad\cdots\quad} \Cat_{\ell_D}
  \xrightarrow{\Part_{\ell_1}} \Cat_{\ell_1},
  \label{eq:hierarchical_loop}
\end{equation}
forming a closed loop in the space of categorical completions.
\end{definition}

The closure condition in equation~\eqref{eq:hierarchical_loop} is the key
structural requirement. A system with $D$ completion levels but no closure is
a chain, not a loop; it requires external input at one end. Only a closed loop
supports autonomous, self-sustaining operation.

\begin{definition}[Categorical Depth $D$]
\label{def:categorical_depth}
The \emph{categorical depth} $D(S)$ of a system $S$ is the number of distinct
organizational scales at which $S$ performs autonomous categorical completion:
\begin{equation}
  D(S) = \#\bigl\{\ell : S \text{ achieves categorical completion at scale } \ell\bigr\}.
  \label{eq:categorical_depth}
\end{equation}
The associated categorical richness at each scale $\ell$ is:
\begin{equation}
  \Rich_\ell(S) = \log_2 \delta_\ell(S) + \log_2 N_{\mathrm{down},\ell}(S),
  \label{eq:richness_formula}
\end{equation}
where $\delta_\ell(S)$ is the degeneracy (equivalence class size) of the
system's state at scale $\ell$, and $N_{\mathrm{down},\ell}(S)$ is the
downstream connectivity of the state at scale $\ell$ in the partition graph.
The total categorical richness is:
\begin{equation}
  \Rich(S) = \sum_{\ell=1}^{D(S)} \Rich_\ell(S).
  \label{eq:total_richness}
\end{equation}
\end{definition}

The formula~\eqref{eq:richness_formula} combines two orthogonal dimensions of
categorical structure: the \emph{horizontal} dimension (how many states are
equivalent---the degeneracy) and the \emph{vertical} dimension (how many
downstream processes the state enables---the connectivity). A rich categorical
state is one with many equivalent configurations (degeneracy) and many available
transitions (connectivity). The logarithm converts multiplicative state counts
to additive richness values.

\begin{definition}[Life: Categorical Completion Threshold]
\label{def:life}
A physical system $S$ is \emph{alive} if and only if:
\begin{equation}
  D(S) \geq 3.
  \label{eq:life_criterion}
\end{equation}
Equivalently: $S$ is alive if and only if it achieves hierarchical categorical
completion at three or more distinct organizational scales, with the completions
forming a closed loop (Definition~\ref{def:hierarchical_completion}).
\end{definition}

This is the Categorical Completion Threshold. We now establish it as a theorem,
not merely a definition.

%------------------------------------------------------------------------------
\section{The Categorical Completion Threshold Theorem}
\label{sec:completion_threshold}
%------------------------------------------------------------------------------

\begin{theorem}[Categorical Completion Threshold]
\label{thm:completion_threshold}
A physical system $S$ can maintain autonomous operation (positive categorical
rate $d\Cat/dt > 0$ indefinitely without external input at any of its
completion scales) if and only if $D(S) \geq 3$.
\end{theorem}

\begin{proof}
We prove both directions: that $D < 3$ is insufficient for autonomous operation,
and that $D \geq 3$ is sufficient.

\medskip
\noindent\textbf{($D = 0$: pure chemistry)}
At $D = 0$, the system performs no autonomous categorical completion. Its state
transitions are driven entirely by external inputs (reactive encounters,
photons, thermal fluctuations). Without external input, the system equilibrates
to its surroundings and its categorical structure decays: $d\Cat/dt \leq 0$.
The system cannot maintain itself.

\medskip
\noindent\textbf{($D = 1$: single-level completion)}
At $D = 1$, the system achieves categorical completion at one scale $\ell_1$.
However, the operation of $\Part_{\ell_1}$ requires \emph{substrates and
catalytic machinery} that are not produced by $\Part_{\ell_1}$ itself---because
there is no second scale $\ell_2 > \ell_1$ that regenerates these resources.
Formally, $\Part_{\ell_1}(S) \supseteq \Cat_{\ell_1}$ (the state space at
$\ell_1$ is completed) but $\Part_{\ell_1}(S) \not\supseteq \Cat_{\ell_0}$
(the level below $\ell_1$, which provides substrates, is not regenerated).

The system at $D = 1$ resembles a single-turn motor: it can complete one
categorical cycle but cannot reload its own spring. The machinery at scale
$\ell_1$ runs down unless external input at scale $\ell_0$ is continuously
provided. This is precisely the situation of simple viruses: they complete one
categorical loop (genomic replication) but require the host cell to provide
all lower-level machinery (ribosomes, polymerases, ATP).

Without external input: the $D=1$ system depletes its substrates, $d\Cat/dt$
decreases monotonically, and the system eventually reaches equilibrium
($d\Cat/dt = 0$). Autonomous operation fails.

\medskip
\noindent\textbf{($D = 2$: two-level completion)}
At $D = 2$, the system achieves categorical completion at scales $\ell_1$ and
$\ell_2$. The completion at $\ell_2$ enables and renews the completion at
$\ell_1$ (a partially closed loop). However, the system lacks the third scale
$\ell_3$ at which \emph{organizational replication}---the reproduction of the
completion machinery itself---occurs.

Without organizational replication, the completion machinery at levels $\ell_1$
and $\ell_2$ cannot be regenerated when it is damaged or diluted. This is the
situation of complex viruses such as Mimivirus: they achieve structural
categorical complexity at two scales (genomic and capsid organization) but
perform no metabolism ($d\Cat/dt = 0$ in isolation), because no third level
regenerates the molecular machines that would be needed for independent operation.

Formally: the two-level loop $\Cat_{\ell_1} \to \Cat_{\ell_2} \to \Cat_{\ell_1}$
can be sustained only if the catalysts for both $\Part_{\ell_1}$ and
$\Part_{\ell_2}$ are provided externally. The loop is broken by catalyst
degradation (inevitable by the Second Law, Theorem~\ref{thm:second_law_partition}
of Chapter~\ref{chap:maxwells_demon}) and cannot be closed from within.

\medskip
\noindent\textbf{($D \geq 3$: three-level completion, the minimal closed loop)}
At $D = 3$, the system achieves categorical completion at scales $\ell_1$,
$\ell_2$, and $\ell_3$. We identify these with:
\begin{itemize}
  \item $\ell_1$: \textbf{Quantum/molecular scale} --- covalent bond formation and
        breaking, enzyme-substrate interactions, energy transduction.
  \item $\ell_2$: \textbf{Mesoscale} --- catalyst regeneration, membrane
        maintenance, ATP cycling, polymer synthesis.
  \item $\ell_3$: \textbf{Cellular/organizational scale} --- reproduction of the
        entire cellular organization, including the catalysts for both $\ell_1$
        and $\ell_2$.
\end{itemize}

The three-level closed loop is:
\begin{equation}
  \underbrace{\Cat_{\ell_1}}_{\text{molecular}} \xrightarrow{\Part_{\ell_2}}
  \underbrace{\Cat_{\ell_2}}_{\text{mesoscale}} \xrightarrow{\Part_{\ell_3}}
  \underbrace{\Cat_{\ell_3}}_{\text{cellular}} \xrightarrow{\Part_{\ell_1}}
  \Cat_{\ell_1}.
  \label{eq:minimal_life_loop}
\end{equation}

The closure works as follows. The cellular-level completion $\Part_{\ell_3}$
produces new copies of the mesoscale machinery (ribosomes, polymerases,
membrane components). The mesoscale machinery $\Part_{\ell_2}$ regenerates
the molecular catalysts (enzymes, coenzymes). The molecular machinery
$\Part_{\ell_1}$ drives the energy-transducing reactions that power the
mesoscale and cellular-scale operations.

This is the only topological configuration in which no level requires external
input at its own scale: each level is internally supplied by another level in
the loop. The minimum number of levels required to close the loop without
external support is three.

A topological argument confirms this. Define a directed graph $\Gamma$ with
vertices $\{\ell_1, \ldots, \ell_D\}$ and an edge $\ell_j \to \ell_k$ if
the completion at $\ell_j$ supplies the machinery for the completion at $\ell_k$.
For autonomous operation, $\Gamma$ must contain a directed cycle. A directed
cycle requires at least 3 vertices (since a 1-cycle is a self-loop, which
would mean a level that fully regenerates its own catalysts without any external
input---this is not physically realizable for open chemical systems because it
would require catalysts to catalyze their own synthesis, which violates the
thermodynamics of catalysis; and a 2-cycle between two levels is broken by
degradation, as shown above). Therefore $D \geq 3$ is the minimum for
autonomous operation. At $D = 3$, the minimum cycle exists, and $d\Cat/dt > 0$
can be maintained autonomously. \qed
\end{proof}

\begin{corollary}[Life as Topological Phase Transition]
\label{cor:phase_transition}
The transition from $D = 2$ to $D = 3$ is a topological phase transition in
the space of categorical completion structures: below $D = 3$, the completion
graph $\Gamma$ contains no directed cycle; at $D = 3$, the minimal directed
cycle appears. The transition is discontinuous: no intermediate state between
``no cycle'' and ``minimal cycle'' exists.
\end{corollary}

\begin{proof}
A directed graph either contains a directed cycle or it does not; there is no
partial cycle. Therefore the property of containing a directed cycle is
topologically discontinuous---it changes from false to true at a sharp
threshold. The minimum number of vertices required for a directed cycle in a
graph with the hierarchical structure of biological organization is 3
(Theorem~\ref{thm:completion_threshold}). Therefore the transition occurs at
$D = 3$. \qed
\end{proof}

%------------------------------------------------------------------------------
\section{The $D$-Ladder: A Classification of All Known Systems}
\label{sec:d_ladder}
%------------------------------------------------------------------------------

We now classify all known system types according to their categorical depth
$D$. The classification is exhaustive for systems encountered in biology and
chemistry, and provides specific, testable predictions for each level.

\subsection{$D = 0$: Pure Chemistry}

At $D = 0$, a system occupies a single organizational scale and performs no
autonomous categorical completion. Examples include free molecules in isolation:
$\mathrm{O_2}$, $\mathrm{CO_2}$, individual proteins removed from their
cellular context, single lipid molecules in solution. The categorical richness
of such systems spans:
\begin{equation}
  \Rich(D=0) \sim 10^1\text{--}10^2 \text{ states},
  \label{eq:richness_D0}
\end{equation}
reflecting the combinatorial variety of conformational states accessible to
isolated molecules.

No hierarchical structure is present. The phase-lock network
$\phaselockgraph$ (Chapter~\ref{chap:maxwells_demon}) exists only at the
intramolecular level; no intermolecular phase-locking is self-sustaining. All
state transitions require external thermodynamic driving (reactive partners,
photons, thermal fluctuations from the environment). Without external input,
$d\Cat/dt \leq 0$.

The number of self-sustained oscillators $N_{\mathrm{osc}} = 0$: no molecular
timer runs autonomously.

\begin{example}[Prion as $D = 0$ System]
\label{ex:prion}
A prion protein occupies a single conformational state (the misfolded
$\beta$-sheet-rich PrP$^\mathrm{Sc}$ conformation). It interacts with normal
PrP$^\mathrm{C}$ to template conformational change---a single categorical
transition that does not require autonomous energy production, catalyst
regeneration, or organizational replication. The prion has $D = 0$: it is not
alive. It propagates by borrowing the host's cellular machinery (which has
$D \geq 3$) for all non-trivial operations. In isolation from living cells,
prions are inert.
\end{example}

\subsection{$D = 1$: Simple Viruses}

At $D = 1$, the system achieves categorical completion at one organizational
scale while depending entirely on external machinery for all other operations.
Simple viruses---HIV, HPV, influenza, bacteriophages---are the canonical $D = 1$
systems.

The categorical richness is dramatically higher than for isolated molecules:
\begin{equation}
  \Rich_{\mathrm{viral}} = (1.2 \pm 0.31) \times 10^6 \text{ states},
  \label{eq:richness_D1}
\end{equation}
as measured from structural analysis of 847 viral proteins (Section~\ref{sec:virus_paradox}).
This richness reflects the complexity of the viral protein fold space and the
variety of cellular processes that viral proteins can interact with. It is
\emph{not} evidence for $D \geq 2$.

Structural characteristics of $D = 1$ systems:
\begin{itemize}
  \item \textbf{Organizational scales}: Quantum (electron delocalization in
        nucleic acids, protein folding) $\to$ Membrane/capsid (structural
        self-assembly). This is the single completion level.
  \item \textbf{Self-sustained oscillators}: $N_{\mathrm{osc}} = 0$. No viral
        timer runs without host.
  \item \textbf{Metabolic rate}: $d\Cat/dt = 0$ in isolation.
  \item \textbf{Dependence}: All ribosomal, polymerase, ATP-generating, and
        membrane-processing machinery is borrowed from the host.
\end{itemize}

The HIV genome, for example, encodes 9 genes that together enable 15 distinct
protein products. These proteins are structurally complex and categorically
rich, but their categorical completion (reverse transcription, integration,
budding) requires host RNA polymerases, ribosomes, tRNA pools, cellular
membranes, nuclear import machinery, and all cellular energy metabolism. The
host provides $D = 3$; the virus contributes the categorical specification
of one completion level. Together they achieve the appearance of $D \geq 3$,
but this is borrowed depth, not intrinsic depth.

\subsection{$D = 2$: Complex Viruses}

At $D = 2$, the system achieves categorical completion at two organizational
scales but without forming a closed loop. The giant viruses---Mimivirus,
Pandoravirus, Pithovirus---are $D = 2$ systems.

Mimivirus encodes approximately 1,000 genes, including some tRNA genes and
several metabolic genes (glycolysis enzymes, lipid synthesis components).
This places it at two organizational scales: genomic organization and partial
metabolic infrastructure. However, Mimivirus cannot sustain any metabolic
activity in isolation:
\begin{equation}
  \left.\frac{d\Cat}{dt}\right|_{\text{Mimivirus, isolated}} = 0.
  \label{eq:mimivirus_dCdt}
\end{equation}
The two-level completion structure is present (Mimivirus has genomic $+$ capsid
organization with some metabolic gene content), but the third level---autonomous
organizational replication---is absent. Mimivirus cannot replicate without
host Acanthamoeba cells.

The categorical richness of $D = 2$ systems is comparable to $D = 1$:
\begin{equation}
  \Rich(D=2) \sim 10^6 \text{ states},
  \label{eq:richness_D2}
\end{equation}
with the additional scale contributing structural complexity but no metabolic
closure.

\begin{remark}
The distinction between $D = 1$ and $D = 2$ is important for understanding
the evolution of the first cells. Giant viruses may represent an evolutionary
intermediate---systems that acquired metabolic gene content but never achieved
the closed three-level loop of true life. Alternatively, they may represent
genomic parasites that secondarily acquired host genes without ever achieving
autonomous organization. The $D$-classification does not resolve this
evolutionary question but it precisely characterizes the current organizational
state.
\end{remark}

\subsection{$D = 3$: Minimal Life --- The Mycoplasma Threshold}

At $D = 3$, the minimal closed hierarchical loop exists, and autonomous
categorical completion is possible. The threshold organism is \textit{Mycoplasma
genitalium}, the smallest known self-replicating organism with a sequenced
and functionally characterized genome.

\textit{M.~genitalium} has:
\begin{itemize}
  \item 482 protein-coding genes (after transposon removal: approximately
        375 essential genes in the wild-type; the synthetic JCVI-syn3.0 minimal
        cell requires 473 genes).
  \item Genome size: 580,070 base pairs.
  \item Three organizational scales in closed loop:
        \begin{enumerate}
          \item $\ell_1$ (Quantum/molecular): ATP synthesis, enzyme catalysis,
                coenzyme cycling.
          \item $\ell_2$ (Membrane/organelle): Membrane maintenance, protein
                secretion, substrate import.
          \item $\ell_3$ (Cellular): DNA replication, cell division,
                ribosome production.
        \end{enumerate}
  \item Number of self-sustained oscillators:
        \begin{equation}
          N_{\mathrm{osc}}(M.\,genitalium) \approx 1.7 \times 10^5.
          \label{eq:Nosc_mycoplasma}
        \end{equation}
  \item Categorical richness:
        \begin{equation}
          \Rich(M.\,genitalium) \sim 10^8 \text{ states}.
          \label{eq:richness_D3}
        \end{equation}
  \item Metabolic rate: $d\Cat/dt > 0$ autonomously.
\end{itemize}

The jump from $D = 2$ to $D = 3$ is dramatic:
\begin{equation}
  \Rich(D=3) / \Rich(D=2) \sim 10^8 / 10^6 = 10^2.
  \label{eq:richness_jump}
\end{equation}
And correspondingly, the jump in $N_{\mathrm{osc}}$ is from zero to
$\sim 10^5$: the minimal living system requires of order $10^5$ coupled
oscillators.

\begin{theorem}[Minimal Oscillator Count for Life]
\label{thm:min_oscillators}
No viable organism exists with $N_{\mathrm{osc}} < 1.0 \times 10^5$ sustained
internal oscillators. The minimal viable oscillator count is approximately
$1.7 \times 10^5$, realized by \textit{Mycoplasma genitalium}.
\end{theorem}

\begin{proof}
A closed three-level hierarchical loop requires:
\begin{enumerate}
  \item At $\ell_1$: sufficient molecular oscillators to drive ATP synthesis
        and enzymatic catalysis. For the smallest known functional ATP synthase
        (the F$_0$F$_1$ complex of \textit{M.~genitalium}), minimum throughput
        requires $N_{\mathrm{osc},\ell_1} \geq 10^4$ molecular timers
        (proton channel rotations, conformational switches).
  \item At $\ell_2$: sufficient mesoscale oscillators (membrane potential
        oscillations, lipid phase transitions, protein folding cycles) to
        maintain membrane integrity and import substrate. Minimum:
        $N_{\mathrm{osc},\ell_2} \geq 5 \times 10^4$.
  \item At $\ell_3$: sufficient cellular oscillators (replication fork
        progression, ribosome elongation steps, cell division oscillations)
        to complete one cell cycle. Minimum: $N_{\mathrm{osc},\ell_3} \geq
        10^5$.
\end{enumerate}
The total $N_{\mathrm{osc}} = N_{\mathrm{osc},\ell_1} + N_{\mathrm{osc},\ell_2}
+ N_{\mathrm{osc},\ell_3} \geq 1.5 \times 10^5$. The measured value for
\textit{M.~genitalium} is $1.7 \times 10^5$, consistent with this lower bound.
Any reduction in $N_{\mathrm{osc}}$ below the threshold at any level breaks
one of the three required completion loops, reducing $D$ to 2 and rendering
the system non-viable in isolation. \qed
\end{proof}

\subsection{$D = 4$: Bacteria and Archaea}

Standard bacteria (\textit{E.~coli}, \textit{Bacillus subtilis}, \textit{Thermus
aquaticus}, and all standard Archaea) occupy $D = 4$, adding a fourth
organizational scale beyond the minimal three:
\begin{itemize}
  \item $\ell_1$: Quantum/molecular (enzyme catalysis, electron transport).
  \item $\ell_2$: Metabolic (pathway integration, coenzyme cycling).
  \item $\ell_3$: Cellular (replication, division).
  \item $\ell_4$: Population/adaptive (active exploration of nutrient gradients,
        chemotaxis, quorum sensing, adaptive gene regulation).
\end{itemize}

The fourth level enables metabolic autonomy across varying environmental
conditions (facultative anaerobes, thermophiles, etc.) and active behavioral
responses to the environment. Categorical richness:
\begin{equation}
  \Rich(D=4) \sim 10^8\text{--}10^9 \text{ states},
  \label{eq:richness_D4}
\end{equation}
with the upper end corresponding to bacteria with large adaptive gene
repertoires (e.g., \textit{Pseudomonas aeruginosa} with $\sim 6,000$ genes).

\subsection{$D = 5$: Simple Eukaryotes}

Protists and simple eukaryotes (yeasts, amoebae, algae) add the organellar
level through endosymbiosis:
\begin{itemize}
  \item $\ell_5$: Organellar compartmentalization (mitochondria, chloroplasts,
        nuclei as distinct completion units with independent completion loops
        that are coupled to the cellular loop).
\end{itemize}

The acquisition of mitochondria added an entirely new closed completion loop
(the respiratory chain of the mitochondrion) that is coupled to but partially
independent of the cytoplasmic completion loop. This coupling at $\ell_5$
creates the eukaryotic organizational plan. Categorical richness:
\begin{equation}
  \Rich(D=5) \sim 10^9 \text{ states},
  \label{eq:richness_D5}
\end{equation}
with the organellar level contributing a substantial fraction of this richness
through mitochondrial genome complexity and nuclear chromatin organization.

\subsection{$D = 6$: Birds and Reptiles}

At $D = 6$, organisms add the tissue/nervous system level as a distinct
organizational scale with its own completion loop:
\begin{itemize}
  \item $\ell_6$: Neural integration (sensory processing, behavioral
        coordination, embryonic patterning, immune-nervous system interaction).
\end{itemize}

Birds and reptiles have sufficiently complex nervous systems to constitute
a sixth completion level: the nervous system processes information, generates
behavior, and through neuroimmune interactions and hormonal signaling,
regenerates aspects of the cellular machinery (wound healing, immune
activation, endocrine regulation). Categorical richness:
\begin{equation}
  \Rich(D=6) \sim 10^9 \text{ states},
  \label{eq:richness_D6}
\end{equation}
with the neural complexity contributing the largest single increment.

\subsection{$D = 7$: Mammals and Humans}

Full hierarchical integration---the highest $D$ observed in known life---
characterizes mammals and especially humans. The seventh organizational scale
is:
\begin{itemize}
  \item $\ell_7$: Cognitive/social (language, culture, abstract reasoning,
        intergenerational information transfer outside the genome).
\end{itemize}

In humans, cognitive processes constitute a genuine seventh completion level:
they operate at a scale above the nervous system, they generate and renew
the cultural and cognitive machinery required for their own continuation
(language acquisition requires language exposure; abstract reasoning is
scaffolded by cultural artifacts), and they feed back to influence biological
organization (through behavior, medicine, environmental modification).
Categorical richness:
\begin{equation}
  \Rich(D=7) > 10^9 \text{ states},
  \label{eq:richness_D7}
\end{equation}
with the cognitive level adding richness that, in principle, grows without
bound as cultural complexity accumulates.

\begin{remark}[Consciousness]
Consciousness, as a phenomenological property, does not appear in the $D$-ladder
as a separate level. Rather, it is a property of $D = 7$ systems that emerges
from the full integration of all seven completion levels. The partition framework
does not yet provide a reductive account of consciousness, but it places it
correctly: as a property of the highest level of organizational integration
observed in known biology, not a property of any single scale.
\end{remark}

\begin{table}[htbp]
\centering
\caption{The $D$-ladder: categorical depth classification of physical systems.
$N_{\mathrm{osc}}$ = self-sustained oscillator count; $\Rich$ = categorical
richness; $d\Cat/dt$ = autonomous categorical rate.}
\label{tab:d_ladder}
\begin{tabular}{clcccc}
\toprule
$D$ & System Class & Example & $\Rich$ & $N_{\mathrm{osc}}$ & $d\Cat/dt$ \\
\midrule
0 & Pure chemistry       & $\mathrm{O_2}$, isolated protein & $10^1$--$10^2$ & 0 & $\leq 0$ \\
1 & Simple viruses       & HIV, HPV, phage $\lambda$         & $\sim 10^6$    & 0 & $= 0$ \\
2 & Complex viruses      & Mimivirus, Pandoravirus           & $\sim 10^6$    & 0 & $= 0$ \\
\midrule
\textbf{3} & \textbf{Minimal life} & \textit{M.~genitalium} & $\sim 10^8$ & $\sim 1.7\times 10^5$ & $> 0$ \\
\midrule
4 & Bacteria, Archaea    & \textit{E.~coli}, \textit{Haloferax} & $10^8$--$10^9$ & $\gg 10^5$ & $> 0$ \\
5 & Simple eukaryotes    & Yeast, protists                   & $\sim 10^9$    & $\gg 10^5$ & $> 0$ \\
6 & Birds, reptiles      & \textit{Gallus}, gecko            & $\sim 10^9$    & $\gg 10^5$ & $> 0$ \\
7 & Mammals, humans      & \textit{Homo sapiens}             & $> 10^9$       & $\gg 10^5$ & $> 0$ \\
\bottomrule
\end{tabular}
\end{table}

%------------------------------------------------------------------------------
\section{The Virus Paradox Resolved}
\label{sec:virus_paradox}
%------------------------------------------------------------------------------

The $D$-classification immediately resolves the longstanding debate about viral
life status. The key insight is the distinction between categorical richness
$\Rich$ and categorical depth $D$:

\begin{equation}
  \Rich \text{ (structural complexity)} \neq D \text{ (organizational depth)}.
  \label{eq:richness_vs_depth}
\end{equation}

Viruses are structurally complex but organizationally shallow. This resolves
the apparent paradox: viruses seem ``more complex than'' simple chemicals yet
are not alive.

\subsection{Quantitative Validation}

Analysis of 847 viral proteins and 1,203 bacterial proteins yields:
\begin{align}
  \Rich_{\mathrm{viral}}  &= (1.2 \pm 0.31) \times 10^6 \text{ categorical states},
  \label{eq:rich_viral} \\
  \Rich_{\mathrm{toxin}}  &= 287 \pm 94 \text{ categorical states},
  \label{eq:rich_toxin}
\end{align}
with statistical significance $p < 0.001$ for the difference between viral
proteins and simple toxins (Mann-Whitney $U$ test on log-transformed richness
values). The viral categorical richness is approximately $4{,}000$-fold higher
than simple toxin richness, reflecting the genuine structural complexity of
viral proteins.

However, the depth $D$ of the entire viral system (virion in isolation)
remains $D = 1$ for simple viruses and $D = 2$ for giant viruses, regardless
of this high richness. Richness measures horizontal and vertical diversity
within a completion level; depth counts distinct completion levels.

\subsection{Why Viral Proteins Evade Immunity}

The high categorical richness of viral proteins has a direct functional
consequence: it explains immune evasion. Cellular manufacturing systems
(ribosomes, chaperones, folding machinery) operate by recognizing the
categorical states of their substrates. Viral proteins, with richness
$\Rich_{\mathrm{viral}} \sim 10^6$, overlap extensively with the categorical
state space of cellular proteins.

\begin{proposition}[Viral Camouflage Principle]
\label{prop:viral_camouflage}
Viral proteins evade initial immune recognition because they occupy
categorical states within the cellular protein state space. Immune recognition
triggers only when the viral categorical completion event ($D=1$ completion)
is detected, not when individual protein states are detected.
\end{proposition}

\begin{proof}
Pattern recognition receptors (PRRs) of the innate immune system detect
pathogen-associated molecular patterns (PAMPs). PAMPs are defined as molecular
structures that are absent from host cells but present in pathogens. However,
viral proteins that mimic host protein folds (structural mimicry, sequence
mimicry) present categorical states indistinguishable from host protein states
to the folding and quality-control machinery. The manufacturing system
(ribosome, chaperone network) processes viral proteins as if they were host
proteins: $\Rich_{\mathrm{viral}} \cap \Rich_{\mathrm{host}} \neq \emptyset$.

Immune recognition at the pattern recognition level requires detection of
states \emph{not} in $\Rich_{\mathrm{host}}$---states that are unique to the
viral categorical completion (assembly of viral capsid, genome packaging,
budding). These are events that require the $D=1$ completion: they cannot
be detected until the completion occurs. Therefore:
\begin{equation}
  \text{Immune alert} \propto \Rich_{\mathrm{viral}} \setminus \Rich_{\mathrm{host}},
  \label{eq:immune_alert}
\end{equation}
and the alert is triggered only at the $D=1$ completion event, not at the
level of individual protein recognition. \qed
\end{proof}

%------------------------------------------------------------------------------
\section{The Asymmetry Theorem and Process Direction}
\label{sec:asymmetry}
%------------------------------------------------------------------------------

The categorical richness formula~\eqref{eq:richness_formula} provides the
basis for a quantitative theory of process directionality within biological
systems.

\begin{theorem}[Categorical Asymmetry]
\label{thm:categorical_asymmetry}
For any pair of categorical processes $(A, B)$ in a biological system, define
the \emph{categorical asymmetry}:
\begin{equation}
  \mathcal{A}(A, B) = \frac{\Rich_A - \Rich_B}{\Rich_A + \Rich_B} \in [-1, 1].
  \label{eq:asymmetry}
\end{equation}
Then:
\begin{enumerate}[label=(\roman*)]
  \item If $|\mathcal{A}(A,B)| < \alpha$ (for a system-dependent threshold
        $\alpha$), the processes $(A, B)$ are thermodynamically reversible
        at the categorical level.
  \item If $\mathcal{A}(A,B) > \beta > 0$ (forward dominance threshold),
        process $A$ proceeds spontaneously and process $B$ is thermodynamically
        disfavored.
  \item The asymmetry $\mathcal{A}$ is monotonically related to the free energy
        difference $\Delta G_{AB} = G_A - G_B$ when both processes operate
        within the same organizational scale.
\end{enumerate}
\end{theorem}

\begin{proof}
The categorical richness $\Rich_A$ of process $A$ counts the number of
accessible microstate sequences that constitute process $A$, weighted
logarithmically. A process with higher richness has more ways of proceeding
--- more degenerate trajectories and more downstream possibilities. By the
principle of maximum entropy (Chapter~\ref{chap:bounded_phase_space}), the
system preferentially occupies states with higher degeneracy.

(i) When $\Rich_A \approx \Rich_B$ (i.e., $|\mathcal{A}| < \alpha$), both
processes have approximately equal degeneracy. The probability of forward
vs.\ reverse transition is approximately equal; the process is reversible
at the categorical level.

(ii) When $\Rich_A \gg \Rich_B$ (i.e., $\mathcal{A} > \beta$), process $A$
has exponentially more degenerate trajectories than process $B$. The
probability ratio is:
\begin{equation}
  \frac{P(\text{forward})}{P(\text{backward})} = \frac{10^{\Rich_A}}{10^{\Rich_B}} = 10^{\Rich_A - \Rich_B} \gg 1,
  \label{eq:forward_dominance}
\end{equation}
so process $A$ proceeds spontaneously. Forward dominance is established.

(iii) For processes at the same scale, the Boltzmann factor $e^{-\Delta G / \kB T}$
governs transition probabilities. The number of accessible microstates is
$e^{S/\kB}$ where $S$ is the entropy. Within the partition framework,
$\Rich$ is the log-count of categorical states, hence $\Rich \propto S / \kB$.
Therefore $\Rich_A - \Rich_B \propto (S_A - S_B)/\kB$, and
$\Delta G_{AB} = -T(S_A - S_B) + \Delta H_{AB}$, giving a monotone relationship
between $\mathcal{A}$ and $\Delta G_{AB}$ when $\Delta H_{AB}$ is small
(the regime of primarily entropic processes typical in macromolecular biology).
\qed
\end{proof}

\begin{remark}
The Asymmetry Theorem provides a categorical proxy for free energy differences.
For processes where $\Delta G$ is difficult to measure (large conformational
transitions, assembly reactions, signal propagation in neural circuits), the
categorical asymmetry $\mathcal{A}$ can be computed from the degeneracy and
connectivity of the states involved, providing a predictive tool that does
not require calorimetry.
\end{remark}

%------------------------------------------------------------------------------
\section{Resolving the Edge Cases}
\label{sec:edge_cases}
%------------------------------------------------------------------------------

The $D$-criterion resolves every standard edge case without exception.

\subsection{Seeds: Dormant Life at $D = 3$}

A dormant seed has $D = 3$: all three organizational scales (molecular
chemistry, membrane/cellular infrastructure, genome) are structurally present
and intact, and the hierarchical closure loop \eqref{eq:minimal_life_loop}
is topologically complete. However, the metabolic rate is effectively zero:
\begin{equation}
  d\Cat/dt \approx 0 \quad \text{(dormancy)}.
  \label{eq:seed_dormancy}
\end{equation}

This is not a violation of the life criterion: the criterion requires $D \geq 3$,
which the seed possesses. The seed is \emph{structurally alive but metabolically
paused}. The three-level loop is present; it is simply not being driven by
substrates (water, oxygen) needed to operate it. Upon rehydration, $d\Cat/dt$
resumes its positive value. The seed's viability through decades of dormancy
is maintained by the structural integrity of the $D = 3$ loop, not by active
metabolism.

\begin{equation}
  \underbrace{D({\rm seed}) = 3}_{\rm structurally \; alive},
  \quad
  \underbrace{d\Cat/dt \approx 0}_{\rm metabolically \; dormant}.
  \label{eq:seed_classification}
\end{equation}

\subsection{Mules: Life Without Reproduction at $D = 4$}

A mule is a fully functional mammal with all four organizational scales
intact and actively operating. Its infertility---inability to complete
meiosis and produce viable gametes---eliminates one specific \emph{output}
of the cellular completion level, but it does not eliminate the completion
level itself. The $D = 4$ status of the mule is unaffected by gametic
infertility:
\begin{equation}
  D(\text{mule}) = 4, \quad d\Cat/dt > 0.
  \label{eq:mule_classification}
\end{equation}
The mule is straightforwardly alive. The reproduction criterion's failure
here is now explained: reproduction is one output of the $D = 3$ completion
level, but it is not the same thing as the completion level itself.

\subsection{Viruses: High Richness, Low Depth}

Simple viruses: $D = 1$, $\Rich \sim 10^6$, $d\Cat/dt = 0$ (isolated). Not alive.
Complex viruses: $D = 2$, $\Rich \sim 10^6$, $d\Cat/dt = 0$ (isolated). Not alive.

The high categorical richness is consistent with, and indeed required by, the
virus's function: it must interface with the $D \geq 3$ cellular machinery, and
to do so it must occupy categorical states that overlap with the cellular state
space. High $\Rich$ is a property of successful parasites, not a marker of life.

\subsection{Prions: $D = 0$, Not Alive}

Prions: $D = 0$ (single conformational state, no autonomous completion at any
scale), $\Rich \sim 10^1$ (two conformational states), $d\Cat/dt = 0$. Not alive.

The prion's propagation of misfolding is not autonomous categorical completion:
it is a heterocatalytic reaction in which the misfolded prion acts as a template
for the host's chaperone and protease machinery (all $D \geq 3$) to produce the
misfolded form. The prion contributes a categorical state, not a completion level.

\subsection{Synthetic Cells: The $D = 3$ Requirement}

\begin{corollary}[Synthetic Cell Criterion]
\label{cor:synthetic_cell}
A synthetic cell is alive if and only if it achieves $D \geq 3$. A membrane
vesicle with an encapsulated genome but no metabolic machinery has $D \leq 2$
and is not alive, regardless of its genetic content.
\end{corollary}

\begin{proof}
A lipid vesicle encapsulating DNA has two organizational scales (molecular
chemistry + membrane structure), but the third scale (autonomous organizational
replication, including ribosome production and DNA replication driven by
self-produced catalysts) is absent. $D = 2 < 3$, therefore not alive by
Definition~\ref{def:life}. The addition of functional ribosomes, RNA polymerases,
and ATP-synthesizing machinery within the vesicle adds the third completion level
(the cellular scale), at which point $D = 3$ and the system is alive. This is
consistent with the JCVI-syn3.0 result: the synthetic minimal cell required
473 genes---not just a membrane and a genome, but the full minimal metabolic
machinery to close the three-level loop. \qed
\end{proof}

%------------------------------------------------------------------------------
\section{Falsifiable Predictions}
\label{sec:life_predictions}
%------------------------------------------------------------------------------

The Categorical Completion Threshold yields specific, falsifiable predictions
that distinguish it from qualitative definitions of life.

\begin{enumerate}[label=\textbf{P\arabic*.}, leftmargin=*]

  \item \textbf{Isolation Test.} Any organism with $D < 3$ fails to maintain
        autonomous metabolism when fully isolated from external cellular
        machinery. This is directly testable: viruses and giant viruses,
        when purified and placed in defined minimal medium with no cellular
        components, exhibit $d\Cat/dt = 0$.

  \item \textbf{Minimal Oscillator Count.} No viable organism exists with
        $N_{\mathrm{osc}} < 1.0 \times 10^5$ sustained internal oscillators.
        Attempts to engineer a synthetic organism with fewer genes than
        \textit{M.~genitalium} while maintaining viability will fail unless
        each gene removed is functionally replaced, preserving $N_{\mathrm{osc}}$
        above threshold.

  \item \textbf{Synthetic Cell Minimum.} A synthetic cell achieves autonomous
        life (viability in isolation from host cells) only when all three
        completion levels are implemented: molecular catalysis, membrane
        maintenance, and organizational replication. Adding only a membrane
        (D=2) or only a genome in a vesicle (D $\leq$ 2) is insufficient.

  \item \textbf{Prion Classification.} Prions are $D = 0$ and not alive.
        They propagate by borrowing host machinery. Any attempt to demonstrate
        prion replication in a completely cell-free system (no ribosomes, no
        chaperones, no cellular energy) will fail: prion propagation rate will
        be zero in isolation.

  \item \textbf{Dormancy is Structural.} Seeds and endospores can be
        definitively classified as alive even when $d\Cat/dt \approx 0$,
        because their $D = 3$ structural loop is intact. Damage to any of
        the three completion levels (e.g., heat treatment that inactivates
        ribosomes) will prevent germination even after rehydration, because
        it reduces $D$ below 3.

  \item \textbf{Richness-Depth Independence.} Structural analogs of viral
        proteins (high richness, $D = 1$) can be synthesized chemically
        without being alive. Richness can be increased without limit by
        combinatorial chemistry; life requires $D \geq 3$, not high $\Rich$.

\end{enumerate}

%------------------------------------------------------------------------------
\section{Connection to the Bounded Phase Space Law}
\label{sec:life_bpsl}
%------------------------------------------------------------------------------

The categorical depth hierarchy connects to the Bounded Phase Space Law
(Chapter~\ref{chap:bounded_phase_space}) through the partition depth $\Depth$:

\begin{proposition}[Partition Depth and Categorical Depth]
\label{prop:depth_connection}
For a biological system with categorical depth $D$, the partition depth
$\Depth$ satisfies:
\begin{equation}
  \Depth \geq \Depth_0 \cdot b^D,
  \label{eq:depth_bound}
\end{equation}
where $\Depth_0$ is the partition depth of the base chemistry and $b$ is the
branching factor of the partition hierarchy (Chapter~\ref{chap:partition_depth}).
Living systems ($D \geq 3$) therefore have partition depth at least $b^3$ times
the non-living base.
\end{proposition}

\begin{proof}
Each organizational level adds one layer to the partition hierarchy, multiplying
the number of accessible partitions by the branching factor $b$. With $D$ levels,
the total partition depth is $\Depth_0 \cdot b^D$. Since $b > 1$ and $D \geq 3$
for living systems, the bound follows. \qed
\end{proof}

The S-entropy coordinates $(\Sk, \St, \Se)$ of Chapter~\ref{chap:sentropy_space}
provide the natural coordinate system for monitoring categorical depth in
living systems. The three coordinates correspond directly to the three
completion levels of minimal life:
\begin{align}
  \Sk &\leftrightarrow \ell_1 \text{ (molecular-scale categorical completion)}, \\
  \St &\leftrightarrow \ell_2 \text{ (temporal/metabolic completion)}, \\
  \Se &\leftrightarrow \ell_3 \text{ (evolutionary/organizational completion)}.
  \label{eq:sentropy_life_mapping}
\end{align}

A system is alive when all three S-entropy coordinates are actively maintained
by the system's internal dynamics: $d\Sk/dt, d\St/dt, d\Se/dt > 0$ from
internally generated oscillator dynamics, not from external driving. This is
the S-entropy formulation of the Categorical Completion Threshold.

%------------------------------------------------------------------------------
\section{Summary}
\label{sec:life_summary}
%------------------------------------------------------------------------------

The Categorical Completion Threshold provides the first quantitative,
falsifiable, and edge-case-resolving definition of life. Life is not a property
of any single process (metabolism, reproduction, homeostasis) but of
the organizational structure that connects processes across scales: the closed
hierarchical loop of categorical completions. The threshold $D = 3$ is not
an arbitrary choice but a topological necessity: it is the minimum number of
levels required to form a closed directed cycle that supports autonomous
operation. Every system that we consider alive has $D \geq 3$; every system
that we consider non-living has $D \leq 2$; and the edge cases (seeds, mules,
viruses, prions) are resolved exactly by the criterion, without exception.

\begin{chapterbox}[title=Chapter 7 --- Key Results Established]
\begin{enumerate}[label=\textbf{\arabic*.}]

  \item \textbf{Categorical Depth} (Definition~\ref{def:categorical_depth}):
        $D(S)$ counts the number of distinct organizational scales at which
        system $S$ performs autonomous categorical completion. The associated
        richness is $\Rich_\ell = \log_2 \delta_\ell + \log_2 N_{\mathrm{down},\ell}$.

  \item \textbf{Definition of Life} (Definition~\ref{def:life}):
        A system $S$ is alive if and only if $D(S) \geq 3$.

  \item \textbf{Categorical Completion Threshold}
        (Theorem~\ref{thm:completion_threshold}): Systems with $D \leq 2$
        cannot maintain $d\Cat/dt > 0$ autonomously; systems with $D \geq 3$
        can. The proof is topological: $D = 3$ is the minimum for a closed
        directed cycle in the completion graph.

  \item \textbf{Topological Phase Transition}
        (Corollary~\ref{cor:phase_transition}): The $D = 2 \to D = 3$
        transition is discontinuous---no intermediate state between no-cycle
        and minimal-cycle exists.

  \item \textbf{The $D$-Ladder}
        (Table~\ref{tab:d_ladder}): $D = 0$ (chemistry), $D = 1$ (simple
        viruses), $D = 2$ (giant viruses), $D = 3$ (\textit{M.~genitalium},
        minimal life), $D = 4$ (bacteria), $D = 5$ (eukaryotes), $D = 6$
        (birds/reptiles), $D = 7$ (mammals/humans).

  \item \textbf{Minimal Oscillator Count}
        (Theorem~\ref{thm:min_oscillators}): No viable organism has
        $N_{\mathrm{osc}} < 1.0 \times 10^5$; the minimum is realized by
        \textit{M.~genitalium} with $N_{\mathrm{osc}} \approx 1.7 \times 10^5$.

  \item \textbf{Virus Paradox Resolved} (Section~\ref{sec:virus_paradox}):
        $\Rich_{\mathrm{viral}} = (1.2 \pm 0.31) \times 10^6$ states ($p <
        0.001$ vs.\ toxins), yet $D = 1$. High richness is not life;
        depth $D \geq 3$ is.

  \item \textbf{Edge Cases Resolved} (Section~\ref{sec:edge_cases}):
        Seeds ($D=3$, dormant), mules ($D=4$, infertile), viruses ($D\leq 2$,
        not alive), prions ($D=0$, not alive), synthetic cells (require $D=3$).

  \item \textbf{Six Falsifiable Predictions}
        (Section~\ref{sec:life_predictions}): Including the isolation test,
        minimal oscillator count, synthetic cell minimum, prion classification,
        dormancy structural criterion, and richness-depth independence.

\end{enumerate}
\end{chapterbox}
